\uuid{0azd}
\exo7id{5403}
\titre{exo7 5403}
\auteur{rouget}
\organisation{exo7}
\datecreate{2010-07-06}
\isIndication{false}
\isCorrection{true}
\chapitre{Continuité, limite et étude de fonctions réelles}
\sousChapitre{Autre}
\cpge{mpsi-mp2i.fonctions-valeurs-reelles-complexes.a, pcsi.techniques-calcul-fonctions.a, ptsi.techniques-calcul-fonctions.a, mpsi-mp2i.nombres-reels-suites.a, pcsi.raisonnement-vocabulaire-ensembliste.c, ptsi.raisonnement-vocabulaire-ensembliste.c, bcpst1.fonctions-reelles-usuelles.a, bcpst1.nombres-reels}
\module{Analyse}
\niveau{L1}
\difficulte{}

\contenu{
\question{
Trouver un exemple de fonction pèriodique dont le groupe des pèriodes est dense dans $\Rr$ mais pas $\Rr$.
}
\reponse{
Soit $f$ la fonction caractéristique de $Q$. Le groupe des périodes de $f$ est $\Qq$. En effet,

$$\forall x\in\Rr,\;\forall r\in\Qq,\;x+r\in\Qq\Leftrightarrow x\in\Qq,$$ et donc 

$\forall x\in\Rr,\;\forall r\in\Qq,\;f(x+r)=f(x)$. Mais on a aussi $$\forall x\in\Rr,\;\forall r\in(\Rr\setminus\Qq),\;x+r\in\Qq,\Leftrightarrow x\notin\Qq,$$ et donc $\forall x\in\Rr,\;\forall r\in(\Rr\setminus\Qq),\;f(x+r)\neq f(x)$.
}
}
